\documentclass[11pt]{article}
\usepackage[margin=1in]{geometry}
\usepackage{amsmath,amssymb,microtype}
\usepackage[hidelinks]{hyperref}

\newcommand{\R}{\mathbb R}

\title{A negative answer to the cusp finite-subset-retraction question\\
in arXiv:2004.09419 via arXiv:2108.03535}
\author{Literature-status note}
\date{13 August 2026}

\begin{document}
\maketitle

\begin{center}
\fbox{\parbox{0.91\textwidth}{\textbf{Status: literature-implied negative
answer.}  The cusp has no Lipschitz retraction $X(4)\to X(3)$, hence it
cannot have the family $X(n)\to X(n-1)$ requested in the source.  This is a
direct identification with a later theorem, not a new result.}}
\end{center}

\section{The source question}

Let
\[
 X=\{(s,|s|^{1/2}):s\in[-1,1]\}\subset\R^2,
\]
with the restricted Euclidean metric.  Question~5.4.1 of
arXiv:2004.09419 (PDF page~131, printed page~122), immediately after
Example~5.4.6, asks:
\begin{quote}
If $X$ is the cusp in Example~5.4.6, do we have Lipschitz retractions
$X(n)\to X(n-1)$?
\end{quote}
Here $X(n)$ is the space of nonempty subsets of $X$ of cardinality at most
$n$, with the Hausdorff metric \cite{Ako20}.

\section{The later obstruction}

Kovalev's Theorem~5.1 in arXiv:2108.03535v2 (statement on PDF page~10)
says that if a metric space $Y$ supports an $L$-Lipschitz retraction
$Y(4)\to Y(3)$ and $[0,1]$ admits an $L$-bi-Lipschitz embedding into $Y$,
then there are $r,M>0$ such that every $p,q\in Y$ with $d(p,q)\le r$ can be
joined in $Y$ by a curve of length at most $M d(p,q)$ \cite{Kov21}.  Thus
one retraction $Y(4)\to Y(3)$ already forces uniform local quasiconvexity
whenever $Y$ contains a bi-Lipschitz interval.

Both facts needed to apply the contrapositive are elementary for the cusp.
First, the positive branch contains a bi-Lipschitz interval: for
\(\gamma(u)=(u^2,u)\), $0\le u\le1$,
\[
 |u-v|\le \|\gamma(u)-\gamma(v)\|
 =|u-v|\sqrt{1+(u+v)^2}\le\sqrt5\,|u-v|.
\]
Second, $X$ is not uniformly locally quasiconvex at its vertex.  Put
\[
 p_t=(-t^2,t),\qquad q_t=(t^2,t),\qquad 0<t\le1.
\]
Their ambient distance is $2t^2$.  Every curve in $X$ joining the two
branches must pass through $(0,0)$, so its length is at least
\[
 \|p_t\|+\|q_t\|=2t\sqrt{1+t^2}.
\]
The ratio of this lower bound to $\|p_t-q_t\|$ is
\(\sqrt{1+t^2}/t\to\infty\).  Therefore no constants $r,M$ as in
Theorem~5.1 can exist.

It follows contrapositively that there is no Lipschitz retraction
$X(4)\to X(3)$.  In particular, the answer to Question~5.4.1 is
\[
 \boxed{\text{No.}}
\]

\section{Provenance and scope}

Kovalev defines the Lipschitz Clustering Property by the existence of all
consecutive retractions $Y(n)\to Y(n-1)$; this is the same property asked
about here.  His paper does not cite or name Akofor's cusp question in the
statement or discussion of Theorem~5.1.  The decisive relation is therefore
an agent-identified implication of a later published theorem, appropriately
recorded as a literature-implied answer.  No assertion is made here about
Akofor's other FSR questions.

\begin{thebibliography}{9}
\bibitem{Ako20}
E.~Akofor,
\emph{Metric Geometry of Finite Subset Spaces},
arXiv:2004.09419 (2020), Question~5.4.1.

\bibitem{Kov21}
L.~V. Kovalev,
\emph{Lipschitz clustering in metric spaces},
arXiv:2108.03535v2 (2021), Theorem~5.1;
J. Geom. Anal. 32 (2022).
\end{thebibliography}

\end{document}
